\documentclass[11pt,letterpaper]{article}
\usepackage[utf8]{inputenc}
\usepackage[T1]{fontenc}
\usepackage[french]{babel}
\usepackage{geometry}
\usepackage{amsmath,amssymb,amsfonts}
\usepackage{booktabs}
\usepackage{array}
\usepackage{etoolbox}

\usepackage{tikz}

\title{Exercice d'algèbre linéaire : La rotation des constellations}
\author{}
\date{}

\begin{document}

\maketitle

\section{Situation}

Considérons deux constellations d'étoiles, la Grande Ourse et la Petite Ourse.
Chacune est composée d'étoiles qu'on dispose sur un plan cartésien où l'origine
pointe vers le pôle nord. Les coordonnées de chaque étoile sont :

\begin{center}
  \begin{tabular}{llc}
    \toprule
    \textbf{Constellation} & \textbf{Étoile}                       & \textbf{Coordonnées $(x, y)$} \\
    \midrule
    \textbf{Grande Ourse}  & Dubhe (coin supérieur droit du bol)   &           $(1, -4)$           \\
                           & Merak (coin inférieur droit du bol)   &           $(0, -6)$           \\
                           & Phecda (coin inférieur gauche du bol) &          $(-3, -5)$           \\
                           & Megrez (coin supérieur gauche du bol) &          $(-3, -3)$           \\
                           & Alioth (1er point de la queue)        &          $(-5, -2)$           \\
                           & Mizar (milieu de la queue)            &          $(-6, -1)$           \\
                           & Alkaid (bout de la queue)             &          $(-7, -2)$           \\
    \addlinespace
    \textbf{Petite Ourse}  & Étoile Polaire (Polaris)              &           $(6, 6)$            \\
                           & Yildun (milieu de la queue)           &           $(5, 7)$            \\
                           & Étoile 9 (base de la queue)           &           $(3, 7)$            \\
                           & Étoile 10 (jonction avec le bol)      &           $(2, 6)$            \\
                           & Kochab (coin supérieur gauche)        &           $(-1, 7)$           \\
                           & Pherkad (coin inférieur gauche)       &           $(-1, 5)$           \\
                           & Étoile 11 (coin inférieur droit)      &           $(1, 4)$            \\
    \bottomrule
  \end{tabular}
\end{center}

Voici la disposition initiale des étoiles dans le ciel :

\begin{center}
  \begin{tikzpicture}[scale=0.45,>=latex]
    % Grille et axes
    \draw[step=1cm,gray!30,very thin] (-11,-11) grid (11,11);
    \draw[->, thick] (-11.5,0) -- (11.5,0) node[right] {$x$};
    \draw[->, thick] (0,-11.5) -- (0,11.5) node[above] {$y$};

    % Coordonnées d'origine Grande Ourse (+3 en x)
    \coordinate (G1) at (1,-4);
    \coordinate (G2) at (0,-6);
    \coordinate (G3) at (-3,-5);
    \coordinate (G4) at (-3,-3);
    \coordinate (G5) at (-5,-2);
    \coordinate (G6) at (-6,-1);
    \coordinate (G7) at (-7,-2);

    % Coordonnées d'origine Petite Ourse (+3 en x)
    \coordinate (P1) at (6,6);
    \coordinate (P2) at (5,7);
    \coordinate (P3) at (3,7);
    \coordinate (P4) at (2,6);
    \coordinate (P5) at (-1,7);
    \coordinate (P6) at (-1,5);
    \coordinate (P7) at (1,4);

    % Tracé Grande Ourse
    \draw[blue, thick] (G1) -- (G2) -- (G3) -- (G4) -- cycle;
    \draw[blue, thick] (G4) -- (G5) -- (G6) -- (G7);
    \foreach \p in {G1,G2,G3,G4,G5,G6,G7} \fill[blue] (\p) circle (4pt);

    % Tracé Petite Ourse (le bol et la queue)
    \draw[blue, thick] (P1) -- (P2) -- (P3) -- (P4);
    \draw[blue, thick] (P4) -- (P5) -- (P6) -- (P7) -- cycle;
    \foreach \p in {P1,P2,P3,P4,P5,P6,P7} \fill[blue] (\p) circle (4pt);

    % Ligne d'alignement pointée
    \draw[dashed, thick, gray] (G2) -- (P1);
    % Ajouter les petits points intermédiaires pour illustrer les "5 fois"
    \foreach \i in {1,2,3,4} { \fill[gray] (1+\i, -4+2*\i) circle (2.5pt); }
      \node[gray, right, align=left] at (3, 0) {\footnotesize Prolonger 5 fois};

    % Quelques étiquettes pour se repérer
    \node[blue, below right] at (G1) {\tiny Dubhe};
    \node[blue, below left] at (G2) {\tiny Merak};
    \node[blue, above right] at (P1) {\tiny Polaris};
    \node[blue, below right] at (G7) {\tiny Grande Ourse};
    \node[blue, above left] at (P5) {\tiny Petite Ourse};
  \end{tikzpicture}
\end{center}

Considérons la Terre qui tourne sur son axe. Faisant une rotation complète par
jour, son angle parcouru dans le temps est donné par la fonction suivante :

\[
  \theta(t) = \frac{2\pi}{24}t
            = \frac{\pi}{12}t \quad (\text{où } t \text{ est exprimé en heures})
\]

Une matrice de rotation est une matrice permettant d'effectuer une rotation
vectorielle dans le plan. Voici la matrice de rotation 2D pour un angle $\theta$
(\emph{theta}) donné :

\[
  R(\theta) = \begin{bmatrix}
                \cos\theta & -\sin\theta \\
                \sin\theta & \cos\theta
              \end{bmatrix}
\]

Par exemple après 6h, la Terre aura fait un quart de tour sur elle-même.
Construisons la matrice de rotation d'un quart de tour:

\[
  R_{1/4} = R(6)
          = \begin{bmatrix}
              \cos \frac{\pi}{12}6 & -\sin \frac{\pi}{12}6 \\
              \sin \frac{\pi}{12}6 & \cos \frac{\pi}{12}6
            \end{bmatrix}
          = \begin{bmatrix}
              \cos \frac{\pi}{2} & -\sin \frac{\pi}{2} \\
              \sin \frac{\pi}{2} & \cos \frac{\pi}{2}
            \end{bmatrix}
          = \begin{bmatrix}
              0  & 1 \\
              -1 & 0
            \end{bmatrix}
\]

\section{Questions}

\begin{enumerate}
  \item Dessinez les constellations telles qu'elles seront dans 6 heures. Voici
        un indice pour trouver la position de l'étoile Polaire après 6 heures:

        \[
          \begin{bmatrix}
            0  & 1 \\
            -1 & 0
          \end{bmatrix} \begin{bmatrix}
                          6 \\
                          6
                        \end{bmatrix}
          = ?
        \]

  \item En combien de temps les étoilent reviennent-elles à la position
        initiale?
  \item Démontrez que $R^4_{1/4} = I$
\end{enumerate}

\newpage

\section*{Annexes}

\subsection*{Degrés et radians}

\begin{tikzpicture}[scale=5.3,cap=round,>=latex]
  % draw the coordinates
  \draw[->] (-1.5cm,0cm) -- (1.5cm,0cm) node[right,fill=white] {$x$};
  \draw[->] (0cm,-1.5cm) -- (0cm,1.5cm) node[above,fill=white] {$y$};

  % draw the unit circle
  \draw[thick] (0cm,0cm) circle(1cm);

  \foreach \x in {0,30,...,360}
  {
    % lines from center to point
    \draw[gray] (0cm,0cm) -- (\x:1cm);
    % dots at each point
    \filldraw[black] (\x:1cm) circle(0.4pt);
    % draw each angle in degrees
    \draw (\x:0.6cm) node[fill=white] {$\x^\circ$};
  }

  % draw each angle in radians
  \foreach \x/\xtext in
    {
      30/\frac{\pi}{6}, 45/\frac{\pi}{4}, 60/\frac{\pi}{3}, 90/\frac{\pi}{2},
      120/\frac{2\pi}{3}, 135/\frac{3\pi}{4}, 150/\frac{5\pi}{6}, 180/\pi,
      210/\frac{7\pi}{6}, 225/\frac{5\pi}{4}, 240/\frac{4\pi}{3},
      270/\frac{3\pi}{2}, 300/\frac{5\pi}{3}, 315/\frac{7\pi}{4},
      330/\frac{11\pi}{6}, 360/2\pi}
    \draw (\x:0.85cm) node[fill=white] {$\xtext$};

  \foreach \x/\xtext/\y in
    {
      % the coordinates for the first quadrant
      30/\frac{\sqrt{3}}{2}/\frac{1}{2},
      45/\frac{\sqrt{2}}{2}/\frac{\sqrt{2}}{2},
      60/\frac{1}{2}/\frac{\sqrt{3}}{2},
      % the coordinates for the second quadrant
      150/-\frac{\sqrt{3}}{2}/\frac{1}{2},
      135/-\frac{\sqrt{2}}{2}/\frac{\sqrt{2}}{2},
      120/-\frac{1}{2}/\frac{\sqrt{3}}{2},
      % the coordinates for the third quadrant
      210/-\frac{\sqrt{3}}{2}/-\frac{1}{2},
      225/-\frac{\sqrt{2}}{2}/-\frac{\sqrt{2}}{2},
      240/-\frac{1}{2}/-\frac{\sqrt{3}}{2},
      % the coordinates for the fourth quadrant
      330/\frac{\sqrt{3}}{2}/-\frac{1}{2},
      315/\frac{\sqrt{2}}{2}/-\frac{\sqrt{2}}{2},
      300/\frac{1}{2}/-\frac{\sqrt{3}}{2}}
    \draw (\x:1.25cm) node[fill=white] {$\left( \xtext,\y \right)$};

  % draw the horizontal and vertical coordinates
  % the placement is better this way
  \draw (-1.25cm,0cm) node[above=1pt] {$(-1,0)$} (1.25cm,0cm) node[above=1pt]
    {$(1,0)$} (0cm,-1.25cm) node[fill=white] {$(0,-1)$}
    (0cm,1.25cm) node[fill=white] {$(0,1)$};
\end{tikzpicture}

\newpage

\section*{Corrigé}

\textbf{1. Calcul de l'angle pour $t = 6$ heures :} Pour $t = 6$, l'angle de
rotation est :
\[
  \theta(6) = \frac{\pi}{12} \times 6
            = \frac{\pi}{2} \text{ radians } (90^\circ)
\]

\textbf{2. Évaluation de la matrice de rotation :} Pour
$\theta = \frac{\pi}{2}$, on a $\cos\left( \frac{\pi}{2} \right) = 0$ et
$\sin\left( \frac{\pi}{2} \right) = 1$. La matrice de rotation devient donc :
\[
  R\left( \frac{\pi}{2} \right) = \begin{bmatrix}
                                    0 & -1 \\
                                    1 & 0
                                  \end{bmatrix}
\]

\textbf{3. Exemple de transformation (pour l'étoile Merak) :} Si l'on prend
l'étoile Merak de la Grande Ourse située en $P = \begin{bmatrix}
                                                   0  \\
                                                   -6
                                                 \end{bmatrix}$, sa nouvelle position après 6 heures est :
\[
  P' = R\left( \frac{\pi}{2} \right) P
     = \begin{bmatrix}
         0 & -1 \\
         1 & 0
       \end{bmatrix} \begin{bmatrix}
                       0  \\
                       -6
                     \end{bmatrix}
     = \begin{bmatrix}
         6 \\
         0
       \end{bmatrix}
\]

En appliquant cette même multiplication matricielle à l'ensemble des sommets,
toute la constellation pivote d'un quart de tour autour de l'origine (le pôle
nord).

\vspace{2em}
\textbf{Représentation visuelle de la transformation :}

\begin{center}
  \begin{tikzpicture}[scale=0.45,>=latex]
    % Grille et axes
    \draw[step=1cm,gray!30,very thin] (-11,-11) grid (11,11);
    \draw[->, thick] (-11.5,0) -- (11.5,0) node[right] {$x$};
    \draw[->, thick] (0,-11.5) -- (0,11.5) node[above] {$y$};

    % ---- Points d'origine (bleu ciel) ----
    \coordinate (G1) at (1,-4);
    \coordinate (G2) at (0,-6);
    \coordinate (G3) at (-3,-5);
    \coordinate (G4) at (-3,-3);
    \coordinate (G5) at (-5,-2);
    \coordinate (G6) at (-6,-1);
    \coordinate (G7) at (-7,-2);

    \coordinate (P1) at (6,6);
    \coordinate (P2) at (5,7);
    \coordinate (P3) at (3,7);
    \coordinate (P4) at (2,6);
    \coordinate (P5) at (-1,7);
    \coordinate (P6) at (-1,5);
    \coordinate (P7) at (1,4);

    % Tracé d'origine
    \draw[blue!40, thick] (G1) -- (G2) -- (G3) -- (G4) -- cycle;
    \draw[blue!40, thick] (G4) -- (G5) -- (G6) -- (G7);
    \foreach \p in {G1,G2,G3,G4,G5,G6,G7} \fill[blue!40] (\p) circle (4pt);

    \draw[blue!40, thick] (P1) -- (P2) -- (P3) -- (P4);
    \draw[blue!40, thick] (P4) -- (P5) -- (P6) -- (P7) -- cycle;
    \foreach \p in {P1,P2,P3,P4,P5,P6,P7} \fill[blue!40] (\p) circle (4pt);

    % ---- Points de destination (rouge) ----
    % Application de la matrice R(90) -> (x,y) devient (-y, x)
    \coordinate (G1_new) at (4,1);
    \coordinate (G2_new) at (6,0);
    \coordinate (G3_new) at (5,-3);
    \coordinate (G4_new) at (3,-3);
    \coordinate (G5_new) at (2,-5);
    \coordinate (G6_new) at (1,-6);
    \coordinate (G7_new) at (2,-7);

    \coordinate (P1_new) at (-6,6);
    \coordinate (P2_new) at (-7,5);
    \coordinate (P3_new) at (-7,3);
    \coordinate (P4_new) at (-6,2);
    \coordinate (P5_new) at (-7,-1);
    \coordinate (P6_new) at (-5,-1);
    \coordinate (P7_new) at (-4,1);

    % Tracé d'arrivée
    \draw[red, thick] (G1_new) -- (G2_new) -- (G3_new) -- (G4_new) -- cycle;
    \draw[red, thick] (G4_new) -- (G5_new) -- (G6_new) -- (G7_new);
    \foreach \p in {G1_new,G2_new,G3_new,G4_new,G5_new,G6_new,G7_new} \fill[red]
      (\p) circle (4pt);

    \draw[red, thick] (P1_new) -- (P2_new) -- (P3_new) -- (P4_new);
    \draw[red, thick] (P4_new) -- (P5_new) -- (P6_new) -- (P7_new) -- cycle;
    \foreach \p in {P1_new,P2_new,P3_new,P4_new,P5_new,P6_new,P7_new} \fill[red]
      (\p) circle (4pt);

    % ---- Flèches de rotation ----
    % Courbure à 30 degrés pour montrer le pivot
    \foreach \i in {1,...,7} {
      \draw[->, dashed, black!60, bend right=30] (G\i) to node[midway,
      fill=white, inner sep=1pt] {\tiny $90^\circ$} (G\i_new);
    }
    \foreach \i in {1,...,7}
    {
      \draw[->, dashed, black!60, bend right=30] (P\i) to node[midway,
      fill=white, inner sep=1pt] {\tiny $90^\circ$} (P\i_new);
    }
  \end{tikzpicture}
\end{center}

\end{document}
